\subsubsection*{Vector form and explicit sigma products}
Equation \eqref{eq:sigma-vector} turns a real vector into an algebra element by contraction against the sigma basis:
\begin{align*}
\mathbf A\cdot\sigv &= A_1\sigma_1+A_2\sigma_2+A_3\sigma_3,\\
\mathbf A\cdot\sigv+\mathbf B\cdot\sigv &= (\mathbf A+\mathbf B)\cdot\sigv.
\end{align*}
With Eqs.\ \eqref{eq:basis} and \eqref{eq:sigma-product},
\begin{align*}
(\mathbf A\cdot\sigv)(\mathbf B\cdot\sigv)
&= \sum_{n,m} A_n B_m \sigma_n \sigma_m \notag\\
&= \sum_n A_n B_n \sigma_n^2 \notag\\
&\quad + \sum_{n<m}(A_n B_m-A_m B_n)\sigma_n \sigma_m \notag\\
&= \mathbf A\cdot\mathbf B \notag\\
&\quad + (A_2B_3-A_3B_2)i\sigma_1 \notag\\
&\quad - (A_1B_3-A_3B_1)i\sigma_2 \notag\\
&\quad + (A_1B_2-A_2B_1)i\sigma_3 \notag\\
&= \mathbf A\cdot\mathbf B \notag\\
&\quad + i(\mathbf A\times\mathbf B)\cdot\sigv.
\end{align*}
Table \ref{tab:components-vector-product} records this component split.

The reversed order flips only the antisymmetric term:
\begin{align*}
(\mathbf B\cdot\sigv)(\mathbf A\cdot\sigv)
&= \mathbf A\cdot\mathbf B-i(\mathbf A\times\mathbf B)\cdot\sigv,\\
\frac{(\mathbf A\cdot\sigv)(\mathbf B\cdot\sigv)+(\mathbf B\cdot\sigv)(\mathbf A\cdot\sigv)}{2}
&= \mathbf A\cdot\mathbf B,\\
\frac{(\mathbf A\cdot\sigv)(\mathbf B\cdot\sigv)-(\mathbf B\cdot\sigv)(\mathbf A\cdot\sigv)}{2}
&= i(\mathbf A\times\mathbf B)\cdot\sigv.
\end{align*}
The same product immediately gives
\begin{align*}
(\mathbf A\cdot\sigv)^2 &= \mathbf A^2,\\
(\mathbf A\cdot\sigv)^{2n} &= (\mathbf A^2)^n,\\
(\mathbf A\cdot\sigv)^{2n+1} &= (\mathbf A^2)^n(\mathbf A\cdot\sigv),\\
(\mathbf A\cdot\sigv)\sigv &= \mathbf A+i\,\mathbf A\times\sigv.
\end{align*}
These are the basic sigma-product identities from which the later paravector formulas are built.

\subsubsection*{Products of real scalars, vectors, and paravectors}
Let
\begin{equation*}
X=a+\mathbf A\cdot\sigv,\qquad
Y=b+\mathbf B\cdot\sigv
\qquad
(a,b\in\R,\ \mathbf A,\mathbf B\in\R^3).
\end{equation*}
Using Eq.\ \eqref{eq:dot-cross},
\begin{align*}
XY
&= (a+\mathbf A\cdot\sigv)(b+\mathbf B\cdot\sigv)\notag\\
&= (ab+\mathbf A\cdot\mathbf B)
+ \bigl(a\mathbf B+b\mathbf A+i(\mathbf A\times\mathbf B)\bigr)\cdot\sigv,\\
YX
&= (ab+\mathbf A\cdot\mathbf B)
+ \bigl(a\mathbf B+b\mathbf A-i(\mathbf A\times\mathbf B)\bigr)\cdot\sigv \notag\\
&= (XY)^{\mathrm R}.
\end{align*}
Table \ref{tab:components-real-paravector-product} gives the same product in component form. The scalar part is symmetric, while the imaginary-vector part changes sign with the order.

Two immediate consequences are
\begin{align*}
[X,Y] &= 2i(\mathbf A\times\mathbf B)\cdot\sigv,\\
XY\in\R\oplus\R^3 &\iff \mathbf A\times\mathbf B=0.
\end{align*}
So two real paravectors commute exactly when their vector parts are parallel.

Setting \(Y=X\) removes the antisymmetric term automatically:
\begin{align*}
X^2
&= a^2+\mathbf A^2+2a\,\mathbf A\cdot\sigv.
\end{align*}
If \(X^2=b+\mathbf B\cdot\sigv\), then matching scalar and vector parts gives
\begin{align*}
b &= a^2+\mathbf A^2,\\
\mathbf B &= 2a\,\mathbf A,
\end{align*}
hence
\begin{align*}
a^2 &= \frac{1}{2}\Bigl(b\pm\sqrt{b^2-\mathbf B^2}\Bigr),\\
\mathbf A^2 &= b-a^2.
\end{align*}
This is the real-paravector square-root relation obtained directly from the product decomposition.

\subsubsection*{Products of complex paravectors}
The master product law \eqref{eq:paravector-product} is the complex extension of the real formulas above. Writing
\begin{equation*}
Z=X+iY=S+\mathbf V\cdot\sigv,\qquad
Z'=X'+iY'=S'+\mathbf V'\cdot\sigv,
\end{equation*}
one obtains every later specialization by setting the unused scalar or vector parts to zero. Table \ref{tab:components-paravector-product} gives the full component decomposition of \(ZZ'\), Table \ref{tab:components-paravector-square} gives \(Z^2\), and Tables \ref{tab:components-zzstar}--\ref{tab:components-zcrevz} record the order-sensitive products built from the conjugate and reverse operators.

The scalar part of a product is cyclic:
\begin{equation*}
\eqtext{scalar}(ZZ')=\eqtext{scalar}(Z'Z).
\end{equation*}
In particular, if \(Z'=MZ^{*\mathrm R}\), then
\begin{align*}
\eqtext{scalar}(Z M Z^{*\mathrm R})
&=\eqtext{scalar}(M Z^{*\mathrm R} Z)\notag\\
&=\eqtext{scalar}(M)\,(Z^{*\mathrm R} Z),
\end{align*}
because \(Z^{*\mathrm R} Z\) is scalar.

The order-sensitive variants fall into three algebraic types:
\begin{align*}
ZZ^* &\eqtext{ is a real scalar plus a complex vector},\\
ZZ^{\mathrm R} &\eqtext{ is a real scalar plus a real vector},\\
ZZ^{\mathrm{R}*} &\eqtext{ is a complex scalar}.
\end{align*}
Those three cases are precisely the ones needed later for quadratic forms, inverses, and transformation formulas.

\noindent The remaining algebra tables collect the recovery formulas, product decompositions, and order-dependent component expansions used in later calculations.
\begin{componenttable}[tab:components-paravector-recovery]{Z}
\componentrow{scalar}{a+ib=(Z+Z^{\mathrm{R}*})/2}{}
\componentrow{vector}{(\mathbf A+i\mathbf B)\cdot\sigv=(Z-Z^{\mathrm{R}*})/2}{}
\componentrow{real part}{a+\mathbf A\cdot\sigv=(Z+Z^{\mathrm R})/2}{}
\componentrow{imaginary part}{b+\mathbf B\cdot\sigv=(Z-Z^{\mathrm R})/(2i)}{}
\componentrow{real scalar}{a=(Z+Z^*+Z^{\mathrm R}+Z^{\mathrm{R}*})/4}{}
\componentrow{imaginary scalar}{b=(Z-Z^*-Z^{\mathrm R}+Z^{\mathrm{R}*})/(4i)}{}
\componentrow{real vector}{\mathbf A\cdot\sigv=(Z-Z^*+Z^{\mathrm R}-Z^{\mathrm{R}*})/4}{}
\componentrow{imaginary vector}{\mathbf B\cdot\sigv=(Z+Z^*-Z^{\mathrm R}-Z^{\mathrm{R}*})/(4i)}{}
\end{componenttable}

\noindent The following two tables collect the full component expansions for the generic product \(ZZ'\) and the special case \(Z^2\).
\begin{componenttable}[tab:components-paravector-product]{Z Z'}
\componentrow{real scalar}{aa'-bb'+\mathbf A\cdot\mathbf A'-\mathbf B\cdot\mathbf B'}{}
\componentrow{imaginary scalar}{ab'+ba'+\mathbf A\cdot\mathbf B'+\mathbf B\cdot\mathbf A'}{}
\componentrow{real vector}{a\mathbf A'-b\mathbf B'+a'\mathbf A-b'\mathbf B-\mathbf A\times\mathbf B'-\mathbf B\times\mathbf A'}{}
\componentrow{imaginary vector}{a\mathbf B'+b\mathbf A'+a'\mathbf B+b'\mathbf A+\mathbf A\times\mathbf A'-\mathbf B\times\mathbf B'}{}
\end{componenttable}

\begin{componenttable}[tab:components-paravector-square]{Z^{2}}
\componentrow{real scalar}{a^{2}-b^{2}+\mathbf A^{2}-\mathbf B^{2}}{}
\componentrow{imaginary scalar}{2ab+2\mathbf A\cdot\mathbf B}{}
\componentrow{real vector}{2(a\mathbf A-b\mathbf B)}{}
\componentrow{imaginary vector}{2(a\mathbf B+b\mathbf A)}{}
\end{componenttable}

\noindent The following two tables make the order dependence of \(ZZ^{*}\) and \(Z^{*}Z\) explicit.
\begin{componenttable}[tab:components-zzstar]{Z Z^{*}}
\componentrow{real scalar}{a^{2}+b^{2}-\mathbf A^{2}-\mathbf B^{2}}{}
\componentrow{imaginary scalar}{0}{}
\componentrow{real vector}{-2\,\mathbf A\times\mathbf B}{\(\leftarrow\) differs from \(Z^{*}Z\)}
\componentrow{imaginary vector}{2(a\mathbf B-b\mathbf A)}{}
\end{componenttable}

\begin{componenttable}[tab:components-zstarz]{Z^{*} Z}
\componentrow{real scalar}{a^{2}+b^{2}-\mathbf A^{2}-\mathbf B^{2}}{}
\componentrow{imaginary scalar}{0}{}
\componentrow{real vector}{2\,\mathbf A\times\mathbf B}{\(\leftarrow\) differs from \(ZZ^{*}\)}
\componentrow{imaginary vector}{2(a\mathbf B-b\mathbf A)}{}
\end{componenttable}

\noindent The following two tables play the same role for products involving the reverse operator.
\begin{componenttable}[tab:components-zzrev]{Z Z^{\mathrm R}}
\componentrow{real scalar}{a^{2}+b^{2}+\mathbf A^{2}+\mathbf B^{2}}{}
\componentrow{imaginary scalar}{0}{}
\componentrow{real vector}{2(a\mathbf A+b\mathbf B+\mathbf A\times\mathbf B)}{\(\leftarrow\) differs from \(Z^{\mathrm R}Z\)}
\componentrow{imaginary vector}{\mathbf 0}{}
\end{componenttable}

\begin{componenttable}[tab:components-zrevz]{Z^{\mathrm R} Z}
\componentrow{real scalar}{a^{2}+b^{2}+\mathbf A^{2}+\mathbf B^{2}}{}
\componentrow{imaginary scalar}{0}{}
\componentrow{real vector}{2(a\mathbf A+b\mathbf B-\mathbf A\times\mathbf B)}{\(\leftarrow\) differs from \(ZZ^{\mathrm R}\)}
\componentrow{imaginary vector}{\mathbf 0}{}
\end{componenttable}

\noindent The following two tables show that the conjugate-reverse products collapse to purely scalar and imaginary-scalar parts.
\begin{componenttable}[tab:components-zzcrev]{Z Z^{\mathrm{R}*}}
\componentrow{real scalar}{a^{2}-b^{2}-\mathbf A^{2}+\mathbf B^{2}}{}
\componentrow{imaginary scalar}{2ab-2\mathbf A\cdot\mathbf B}{}
\componentrow{real vector}{\mathbf 0}{}
\componentrow{imaginary vector}{\mathbf 0}{}
\end{componenttable}

\begin{componenttable}[tab:components-zcrevz]{Z^{\mathrm{R}*} Z}
\componentrow{real scalar}{a^{2}-b^{2}-\mathbf A^{2}+\mathbf B^{2}}{}
\componentrow{imaginary scalar}{2ab-2\mathbf A\cdot\mathbf B}{}
\componentrow{real vector}{\mathbf 0}{}
\componentrow{imaginary vector}{\mathbf 0}{}
\end{componenttable}


\subsubsection*{Order dependence and commutators}
The part of a product that depends on order is its commutator. For generic complex paravectors \(Z,Z'\in\C\oplus\C^3\),
\begin{equation}
[Z,Z']
:= ZZ'-Z'Z
= 2i(\mathbf V\times\mathbf V')\cdot\sigv
\label{eq:paravector-commutator}
\end{equation}
Equation \eqref{eq:paravector-commutator} is the compact commutator law. In real and imaginary vector components it expands to
\begin{align*}
[Z,Z']
&= -2(\mathbf A\times\mathbf B'+\mathbf B\times\mathbf A')\cdot\sigv \notag\\
&\quad +2i(\mathbf A\times\mathbf A'-\mathbf B\times\mathbf B')\cdot\sigv.
\end{align*}
The real-paravector commutator
\begin{equation*}
[X,Y]=2i(\mathbf A\times\mathbf B)\cdot\sigv
\end{equation*}
is the corresponding specialization. In every case the scalar terms cancel, so the commutator isolates the cross-product part of the multiplication law.

The same order dependence appears already in the sigma basis itself. The pseudoscalar commutes with the generators,
\begin{align*}
[i,\sigma_n] &= 0,\\
i\sigv &= \sigv i = \{\sigma_2\sigma_3,\ \sigma_3\sigma_1,\ \sigma_1\sigma_2\},
\end{align*}
while distinct sigma generators satisfy
\begin{equation*}
[\sigma_n,\sigma_m]=2\sigma_n\sigma_m=2\epsilon_{nmk}i\sigma_k
\qquad (n\neq m).
\end{equation*}
Thus the oriented sigma-number-two basis is exactly the commutator closure of the sigma-number-one generators.

\subsubsection*{Paravector metric and quadratic forms}
The quadratic forms \eqref{eq:euclidean-form} and \eqref{eq:minkowski-form} are product formulas rather than added structure. For a real paravector \(X=a+\mathbf A\cdot\sigv\),
\begin{align*}
\langle\!\langle X\rangle\!\rangle^2
&= a^2-\mathbf A^2\\
&= XX^*=X^*X.
\end{align*}
The indefinite form is therefore the product of a real paravector with its conjugate.

For a complex paravector \(Z=S+\mathbf V\cdot\sigv=X+iY\),
\begin{align*}
\|Z\|^2
&= |S|^2+\mathbf V\cdot\mathbf V^*\\
&= \eqtext{scalar}(Z^{\mathrm R}Z)
= \eqtext{scalar}(ZZ^{\mathrm R}),
\end{align*}
while the indefinite form is
\begin{align*}
\langle\!\langle Z\rangle\!\rangle^2
&= S^2-\mathbf V^2\\
&= ZZ^{\mathrm{R}*}
= Z^{\mathrm{R}*}Z \notag\\
&= (a^2-b^2-\mathbf A^2+\mathbf B^2)
+ 2i(ab-\mathbf A\cdot\mathbf B).
\end{align*}
Tables \ref{tab:components-zzrev}--\ref{tab:components-zcrevz} show these same product types in component form.

The same formulas also show how the two quadratic forms scale:
\begin{align*}
\|r e^{i\theta} Z\|^2 &= r^2\|Z\|^2,\\
\langle\!\langle r e^{i\theta} Z\rangle\!\rangle^2
&= r^2 e^{2i\theta}\langle\!\langle Z\rangle\!\rangle^2.
\end{align*}
The Euclidean form is insensitive to phase, while the indefinite form retains the phase square.

\subsubsection*{Inverse formulas}
The inverse formulas are the algebraic consequence of the quadratic forms. If \(\mathbf A\times\mathbf B=0\), then
\begin{equation*}
\langle\!\langle XY\rangle\!\rangle^2
= \langle\!\langle X\rangle\!\rangle^2\langle\!\langle Y\rangle\!\rangle^2.
\end{equation*}
In particular, when \(a^2-\mathbf A^2\neq 0\),
\begin{align*}
X^{-1}
&= \frac{X^*}{XX^*}
= \frac{a-\mathbf A\cdot\sigv}{a^2-\mathbf A^2}.
\end{align*}
When \(\langle\!\langle Z\rangle\!\rangle^2\neq 0\), the same argument gives
\begin{align}
Z^{-1}
&= \frac{\crev{Z}}{\langle\!\langle Z\rangle\!\rangle^2}
= \frac{\crev{Z}}{\crev{Z}Z}.
\label{eq:paravector-inverse}
\end{align}
If the corresponding quadratic form vanishes, the inverse does not exist.

Differentiating \(X^{-1}X=1\) gives the standard inverse-product rule
\begin{align*}
\dd(X^{-1}) &= -X^{-1}(\dd X)X^{-1},
\end{align*}
and when \(\mathbf A\times\mathbf B=0\) one also has
\begin{equation*}
X^{-1}Y=Y X^{-1}.
\end{equation*}
